
    %%%%%%%%%%%%%%%%%%%%%%%%%%%%%%%%%%%%%%%%%%%%%%%%%%%%%%%%%%%%%%%%%%%%%%%%%%%%%%%%
    %%% Classe do documento
    \documentclass[12pt, addpoints]{exam}
    \UseRawInputEncoding

    %%%%%%%%%%%%%%%%%%%%%%%%%%%%%%%%%%%%%%%%%%%%%%%%%%%%%%%%%%%%%%%%%%%%%%%%%%%%%%%%
    % preambulo.tex
    %
    % Arquivo de configuração de packages para uso o LaTeX, conforme minhas
    % preferências e modelos pessoais.
    %
    % Para maiores informações, visite:
    %    https://github.com/abrantesasf/latex
    %
    % NÃO ALTERE SE NÃO SOUBER O QUE ESTÁ FAZENDO!
    %%%%%%%%%%%%%%%%%%%%%%%%%%%%%%%%%%%%%%%%%%%%%%%%%%%%%%%%%%%%%%%%%%%%%%%%%%%%%%%%


    %%%%%%%%%%%%%%%%%%%%%%%%%%%%%%%%%%%%%%%%%%%%%%%%%%%%%%%%%%%%%%%%%%%%%%%%%%%%%%%%
    %%% Estruturas de controle:
    \input{static/utils/controles}

    %%%%%%%%%%%%%%%%%%%%%%%%%%%%%%%%%%%%%%%%%%%%%%%%%%%%%%%%%%%%%%%%%%%%%%%%%%%%%%%%
    %%% Compilação condicional em PDF:
    \input{static/utils/pdf}

    %%%%%%%%%%%%%%%%%%%%%%%%%%%%%%%%%%%%%%%%%%%%%%%%%%%%%%%%%%%%%%%%%%%%%%%%%%%%%%%%
    %%% Configurações de layout da página:
    \input{static/utils/layout}

    %%%%%%%%%%%%%%%%%%%%%%%%%%%%%%%%%%%%%%%%%%%%%%%%%%%%%%%%%%%%%%%%%%%%%%%%%%%%%%%%
    %%% Configurações para autores:
    \input{static/utils/autores}

    %%%%%%%%%%%%%%%%%%%%%%%%%%%%%%%%%%%%%%%%%%%%%%%%%%%%%%%%%%%%%%%%%%%%%%%%%%%%%%%%
    %%% Cabeçalho e rodapé:
    %%% Atenção! É MUITO DIFÍCIL ajustar apropriadamente os running headers
    %%% e running footers da primeira vez. Você pode confiar no LaTeX e deixar que
    %%% ele faça o ajuste automaticamente ou pode quebrar a cabeça e
    %%% tentar melhorar algo que o LaTeX já faz muito bem, mesmo que não fique
    %%% exatamente do jeito que você quer. O que você prefere? Se quiser ajustar
    %%% manualmente, descomente a linha abaixo e faça as configurações no arquivo
    %%% "utils/headings.tex".
    %\input{static/utils/headings}

    %%%%%%%%%%%%%%%%%%%%%%%%%%%%%%%%%%%%%%%%%%%%%%%%%%%%%%%%%%%%%%%%%%%%%%%%%%%%%%%%
    %%% Configuração para as fontes do documento:
    %%% Se você quiser usar fontes próprias ou do sistema, precisará ajustar as
    %%% configurações no arquivo "utils/fontes.tex".
    %%% ATENÇÃO: por padrão eu utilizo XeLaTeX com fontes proprietárias projetadas
    %%% por Matthew Butterick, adquiridas em https://mbtype.com, que não podem ser
    %%% distribuídas devido à licença de utilização. Se você usar XeLaTeX, você
    %%% DEVERÁ OBRIGATORIAMENTE ajustar as fontes no arquivo "utils/fontes.tex".
    \input{static/utils/fontes}

    %%%%%%%%%%%%%%%%%%%%%%%%%%%%%%%%%%%%%%%%%%%%%%%%%%%%%%%%%%%%%%%%%%%%%%%%%%%%%%%%
    %%% Sumário:
    \input{static/utils/toc}

    %%%%%%%%%%%%%%%%%%%%%%%%%%%%%%%%%%%%%%%%%%%%%%%%%%%%%%%%%%%%%%%%%%%%%%%%%%%%%%%%
    %%% Matemática:
    \input{static/utils/matematica}

    %%%%%%%%%%%%%%%%%%%%%%%%%%%%%%%%%%%%%%%%%%%%%%%%%%%%%%%%%%%%%%%%%%%%%%%%%%%%%%%%
    %%% Notas de rodapé:
    \input{static/utils/footnotes}

    %%%%%%%%%%%%%%%%%%%%%%%%%%%%%%%%%%%%%%%%%%%%%%%%%%%%%%%%%%%%%%%%%%%%%%%%%%%%%%%%
    %%% Referências cruzadas, links, citações:
    \input{static/utils/crossrefs}

    %%%%%%%%%%%%%%%%%%%%%%%%%%%%%%%%%%%%%%%%%%%%%%%%%%%%%%%%%%%%%%%%%%%%%%%%%%%%%%%%
    %%% Configurações para os hiperlinks em PDF:
    \input{static/utils/pdfconfig}

    %%%%%%%%%%%%%%%%%%%%%%%%%%%%%%%%%%%%%%%%%%%%%%%%%%%%%%%%%%%%%%%%%%%%%%%%%%%%%%%%
    %%% Glossário e índice remissivo:
    \input{static/utils/glossindex}

    %%%%%%%%%%%%%%%%%%%%%%%%%%%%%%%%%%%%%%%%%%%%%%%%%%%%%%%%%%%%%%%%%%%%%%%%%%%%%%%%
    %%% Referências bibliográficas:
    \input{static/utils/bibliografia}

    %%%%%%%%%%%%%%%%%%%%%%%%%%%%%%%%%%%%%%%%%%%%%%%%%%%%%%%%%%%%%%%%%%%%%%%%%%%%%%%%
    %%% Cores:
    \input{static/utils/cores}

    %%%%%%%%%%%%%%%%%%%%%%%%%%%%%%%%%%%%%%%%%%%%%%%%%%%%%%%%%%%%%%%%%%%%%%%%%%%%%%%%
    %%% Computação:
    \input{static/utils/computacao}

    %%%%%%%%%%%%%%%%%%%%%%%%%%%%%%%%%%%%%%%%%%%%%%%%%%%%%%%%%%%%%%%%%%%%%%%%%%%%%%%%
    %%% Gráficos:
    \input{static/utils/graficos}

    %%%%%%%%%%%%%%%%%%%%%%%%%%%%%%%%%%%%%%%%%%%%%%%%%%%%%%%%%%%%%%%%%%%%%%%%%%%%%%%%
    %%% Ícones:
    \input{static/utils/icones}

    %%%%%%%%%%%%%%%%%%%%%%%%%%%%%%%%%%%%%%%%%%%%%%%%%%%%%%%%%%%%%%%%%%%%%%%%%%%%%%%%
    %%% Blocos:
    \input{static/utils/blocos}

    %%%%%%%%%%%%%%%%%%%%%%%%%%%%%%%%%%%%%%%%%%%%%%%%%%%%%%%%%%%%%%%%%%%%%%%%%%%%%%%%
    %%% Boxes coloridos:
    \input{static/utils/colorbox}

    %%%%%%%%%%%%%%%%%%%%%%%%%%%%%%%%%%%%%%%%%%%%%%%%%%%%%%%%%%%%%%%%%%%%%%%%%%%%%%%%
    %%% Tabelas:
    \input{static/utils/tabelas}

    %%%%%%%%%%%%%%%%%%%%%%%%%%%%%%%%%%%%%%%%%%%%%%%%%%%%%%%%%%%%%%%%%%%%%%%%%%%%%%%%
    %%% Ambientes de listas:
    \input{static/utils/listas}

    %%%%%%%%%%%%%%%%%%%%%%%%%%%%%%%%%%%%%%%%%%%%%%%%%%%%%%%%%%%%%%%%%%%%%%%%%%%%%%%%
    %%% Ambientes floats e similares:
    \input{static/utils/floats}

    %%%%%%%%%%%%%%%%%%%%%%%%%%%%%%%%%%%%%%%%%%%%%%%%%%%%%%%%%%%%%%%%%%%%%%%%%%%%%%%%
    %%% Ajustes para a classe EXAM:
    \input{static/utils/exam}

    %%%%%%%%%%%%%%%%%%%%%%%%%%%%%%%%%%%%%%%%%%%%%%%%%%%%%%%%%%%%%%%%%%%%%%%%%%%%%%%%
    %%% Ajustes para textos:
    \input{static/utils/texto}

    %%%%%%%%%%%%%%%%%%%%%%%%%%%%%%%%%%%%%%%%%%%%%%%%%%%%%%%%%%%%%%%%%%%%%%%%%%%%%%%%
    %%% Ambientes, aliases, comandos e customizações em geral para serem
    %%% utilizadas nos documentos. Defina aqui qualquer customização específica
    %%% para seus documentos!
    \input{static/utils/customizacoes}

    %%%%%%%%%%%%%%%%%%%%%%%%%%%%%%%%%%%%%%%%%%%%%%%%%%%%%%%%%%%%%%%%%%%%%%%%%%%%%%%%
    %%% Ajuste do layout, espaçamento de linhas e etc.:
    \geometry{a4paper, portrait, left=2.0cm, right=2.0cm, top=2.0cm, bottom=2.0cm}
    \singlespacing

    %%%%%%%%%%%%%%%%%%%%%%%%%%%%%%%%%%%%%%%%%%%%%%%%%%%%%%%%%%%%%%%%%%%%%%%%%%%%%%%%
    %%% Configurações de cabeçalho e rodapé (não use fancyhdr pois ocorre conflito):
    \pagestyle{headandfoot}
    \runningheadrule
    \firstpageheader{}{}{}
    \runningheader{Sem disciplina}{}{Avaliação}
    \firstpagefooter{}{}{Página \thepage\ de \numpages}
    \runningfooter{}{}{Página \thepage\ de \numpages}
    \runningfootrule

    %%%%%%%%%%%%%%%%%%%%%%%%%%%%%%%%%%%%%%%%%%%%%%%%%%%%%%%%%%%%%%%%%%%%%%%%%%%%%%%%
    %%% Configurações para as propriedades do PDF:
    \hypersetup{
    hidelinks,           % Comente para web, descomente para impressão
    colorlinks=true,      % True para web, False para impressão
    pdftitle=pedro: Avaliação,
    pdfauthor=Matheus Endlich Silveira,
    pdfsubject=Sem disciplina,
    pdfkeywords=['Sem tag'],
    pdfinfo={
        CreationDate={}, % Ex.: D:AAAAMMDDHH24MISS
        ModDate={}       % Ex.: D:AAAAMMDDHH24MISS
    }
    }
    \input{static/utils/pdfconfig}

    %%%%%%%%%%%%%%%%%%%%%%%%%%%%%%%%%%%%%%%%%%%%%%%%%%%%%%%%%%%%%%%%%%%%%%%%%%%%%%%%
    %%% Começa o documento
    \begin{document}

    %%%%%%%%%%%%%%%%%%%%%%%%%%%%%%%%%%%%%%%%%%%%%%%%%%%%%%%%%%%%%%%%%%%%%%%%%%%%%%%%
    %%% Folha de instruções padronizadas da UVV para a prova
    \pdfbookmark[1]{Instruções}{instrucoes}
    \begin{figure}[H]
        \begin{center}
            \includegraphics[scale=0.3]{{static/imgs/logo-uvv.png}}
        \end{center}	
    \end{figure}

    \vspace{-1cm}

    \begin{table}[H]
        \centering
        \begin{tabular}{|llll|}
            \hline
            \multicolumn{2}{|l|}{\textbf{Disciplina:} Sem disciplina} & 
            \multicolumn{1}{l|}{\multirow{3}{*}{\textbf{\makecell{Nota:\\ \,\\ \,}}}} & 
            \multirow{3}{*}{\textbf{\makecell{Coordenador:\\ \,\\ \,}}} \\ 
            \cline{1-2}
            \multicolumn{2}{|l|}{\textbf{Professor:} Matheus Endlich Silveira \hspace{4.72cm} } & 
            \multicolumn{1}{l|}{} & \\ 
            \cline{1-2}
            \multicolumn{2}{|l|}{\textbf{Aluno:}} & 
            \multicolumn{1}{l|}{} & \\ 
            \hline
            \multicolumn{1}{|l|}{\textbf{Turma:} \hspace{6.3cm} } & 
            \multicolumn{1}{l|}{\textbf{Semestre:}} & 
            \multicolumn{2}{l|}{\textbf{Valor:} 10 pontos} \\ 
            \hline
            \multicolumn{1}{|l|}{\textbf{Data:}} & 
            \multicolumn{3}{l|}{\textbf{Avaliação:}} \\ 
            \hline
        \end{tabular}
    \end{table}

    \begin{center}
        \fbox{\setlength\fboxsep{0.3cm} \fbox{\parbox{15.4cm}{
            \begin{center}
                \textbf{INSTRUÇÕES PARA A REALIZAÇÃO DESTA AVALIAÇÃO:}
            \end{center}
            \begin{itemize}[leftmargin=*]
                \item Esta é a primeira avaliação da disciplina Sem disciplina, 
                    e engloba o conteúdo referente Sem tag.
                \item Esta avaliação contém 50 questões objetivas, \textbf{todas obrigatórias}.
                \item \textbf{Leia atentamente cada questão!} As questões objetivas terão uma,
                    e apenas uma única, resposta a ser marcada.
                \item \textbf{Desligue o celular} antes de começar. A avaliação é
                    \textbf{individual} e \textbf{sem consulta}.
                \item As questões podem ser respondidas, na prova, com lápis ou caneta. Ao final
                    da prova \textbf{{VOCÊ DEVERÁ PREENCHER O GABARITO COM CANETA
                    DE COR PRETA}} \textbf{\underline{OU AZUL}} para que seja feita a correção
                    automatizada através da leitura do padrão de respostas marcado no
                    gabarito.
                \item \textbf{CUIDADO ao preencher o gabarito} pois eles são \textbf{INDIVIDUAIS
                    e NOMEADOS} (cada estudante tem um gabarito próprio específico, com um
                    código de identificação único). \textbf{Questões rasuradas são anuladas}
                    pelo software de leitura do gabarito.
                \item Ao final da prova, \textbf{DEVOLVA A PROVA E O GABARITO} para o professor.
                \item Siga todas as normas de \textbf{Integridade Acadêmica} da disciplina.
                    Alunos que forem flagrados com qualquer espécie de ``cola'' ou trocando
                    informações com outros alunos terão suas avaliações recolhidas, as notas
                    zeradas, e a situação será encaminhada para a coordenação para a aplicação
                    das penalidades previstas pela universidade.
                \item Com 90 minutos de avaliação você terá tempo de até 2,min por questão,
                    em média.
                \item Boa avaliação!
            \end{itemize}
            \vspace{2cm}
        }}}
    \end{center}
    \newpage

    %%%%%%%%%%%%%%%%%%%%%%%%%%%%%%%%%%%%%%%%%%%%%%%%%%%%%%%%%%%%%%%%%%%%%%%%%%%%%%%%
    %%% Inicia o bloco de questões:
    \begin{questions}

    \newpage
    \question

Considere a seguinte tabela para uma base de dados relacional:\\

Empregado (CodEmp, NomeEmp, CodDepto)\\

Considere que esta tabela tem um índice na forma de uma árvore B sobre as colunas (CodEmp, CodDepto), nesta ordem.

Quanto a este índice, considere as seguintes afirmativas:

\begin{enumerate}

    \item Este índice pode ser usado pelo SGBD relacional para acelerar uma consulta na qual são fornecidos os valores de CodEmp e CodDepto.

    \item Este índice pode ser usado pelo SGBD relacional para acelerar uma consulta na qual é fornecido um valor de CodEmp.

    \item Este índice não é adequado para ser usado pelo SGBD relacional para acelerar uma consulta na qual é fornecido um valor de CodDepto.

    \item O algoritmo que faz inserções e remoções de entradas do índice tem por objetivo garantir que o índice fique organizado de tal forma que o acesso a cada nodo da árvore implique em número de acessos semelhantes.

    \item O índice por árvore-B não é adequado para tabelas que sofrem grande número de inclusões e exclusões, pois exige reorganizações frequentes.

\end{enumerate}





Quanto a estas afirmativas pode se dizer que:
\begin{choices}
\choice Todas afirmativas estão corretas
\choice Apenas as afirmativas 1), 2) e 5) estão corretas
\choice Apenas as afirmativas 1), 2) e 4) estão corretas
\choice Apenas as afirmativas 1), 2), 3) e 4) estão corretas
\choice Nenhuma das afirmativas está correta
\end{choices}

\question

Para cada \( n \in \mathbb{N} \), seja \( D_n = (0, 1/n) \), onde \( (0, 1/n) \) representa o intervalo aberto de extremos \( 0 \) e \( 1/n \). O conjunto diferença \( D_3 - D_{20} \) é igual a:


\begin{choices}
\choice \( D_{20} \cup D_3 \)
\choice \( [1/20, 1/3) \)
\choice \( D_3 \)
\choice \( (1/20, 1/3) \)
\choice \( D_{20} \)
\end{choices}

\question

A medida da interconexão entre os módulos de uma estrutura de software é denominada e que também é usada em projetos orientados a objetos é: 
\begin{choices}
\choice abstração procedimental
\choice coesão
\choice unidade funcional
\choice ocultamento da informação
\choice acoplamento
\end{choices}

\question

(Adaptado do ENADE 2005) Analise o relacionamento existente entre as tabelas

``clubes' e ``jogadores', nas alternativas abaixo, e marque a alternativa que

indica que todo jogador deve pertencer a um único clube, e todo clube poder ter

ou não jogadores.
\begin{choices}
\choice \includegraphics[width=0.5\linewidth]{static/imgs/defaul.png} \vspace{0.5cm}
\choice \includegraphics[width=0.5\linewidth]{static/imgs/defaul.png} \vspace{0.5cm}
\choice \includegraphics[width=0.5\linewidth]{static/imgs/defaul.png} \vspace{0.5cm}
\choice \includegraphics[width=0.5\linewidth]{static/imgs/defaul.png} \vspace{0.5cm}
\choice \includegraphics[width=0.5\linewidth]{static/imgs/defaul.png} \vspace{0.5cm}
\end{choices}

\question

Um banco de dados OLAP (Online Analytical Processing --- Processamento Analítico

Online) é um banco de dados organizado de forma a dar suporte ao Business

Intelligence (BI --- Inteligência de Negócio) que, de forma simplificada, é o

processo de analisar dados para obter informações que podem ser utilizadas na

tomada de decisões comerciais e estratégicas. Em relação aos bancos de dados

OLAP, marque a afirmação correta:
\begin{choices}
\choice São otimizados para o processamento de grande volume de transações diárias, o que permite a análise e a tomada de decisões em tempo real, favorecendo assim a administração empresarial.
\choice São preparados para responder perguntas do tipo ``Quais alunos estão matriculados na turma de Bancos de Dados I neste semestre?'
\choice São otimizados para o processamento de grande volume de dados históricos e métricas de negócio para a tomada de decisões estratégicas, permitindo e facilitando a análise de dados para a geração de informações e conhecimentos para a tomada de decisões.
\choice Geralmente os dados de origem do OLTP são de bancos de dados OLAP.
\choice São otimizados para consulta e relatórios de rotina de dados.
\end{choices}

\question

Dentre as características do modelo relacional e do modelo de objetos em bancos de dados, qual afirmação é \textbf{INCORRETA}?


\begin{choices}
\choice O modelo de objetos possui mais recursos estruturais para a representação de dados que o relacional.
\choice O modelo de objetos é mais adequado para a representação de tipos abstratos de dados.
\choice O modelo de objetos provê uma representação bem próxima de linguagens de programação.
\choice O relacionamento de herança é diretamente representado no modelo relacional.
\choice O relacionamento binário \( N \times M \) é representado de modo semelhante nos dois modelos.	
\end{choices}

\question

Uma definição lógica, abstrata e auto-contida dos objetos que modelam a

estrutura de armazenamento de dados, dos operadores que modelam o comportamento

e dos demais conceitos que modelam a integridade, e que, juntos, constituem a

máquina abstrata com a qual os usuários interagem é chamada de:


\begin{choices}
\choice Nenhuma das respostas acima
\choice Modelo relacional
\choice Modelo entidade-relacionamento
\choice Banco de dados
\choice Modelo de dados
\end{choices}

\question

Avalie o banco de dados de professores e alunos, mostrado a seguir. Esse banco

de dados pretende armazenar todas as turmas, sendo que um professor pode dar

aula para vários alunos ao mesmo tempo, e um mesmo aluno pode ter aula com

diferentes professores ao mesmo tempo.

\begin{figure}[H]

  \centering

  \caption{Professores, turmas e alunos}

  \label{fig:professores}

  %\fbox{

    \includegraphics[width=0.5\linewidth]{static/imgs/defaul.png}

  %}\\

  %\footnotesize{Fonte: xxxx.}

\end{figure}

Esse banco de dados resolveráo problema de armazenar os dados das turmas?
\begin{choices}
\choice Não, pois as cardinalidades e os sentidos dos relacionamentos entre 'professores' e 'turmas', e entre 'turmas' e 'alunos' não cria um relacionamento N:N entre 'professores' e 'alunos'.
\choice Não é possível determinar somente avaliando o diagrama exibido, seriam necessárias mais informações do dicionário de dados.
\choice Sim, pois o diagrama mostra claramente que um professor pode ter zero ou mais turmas, e cada aluno pode participar de zero ou mais turmas também.
\choice Não, pois o relacionamento entre 'turmas' e 'professores' está errado: a PK 'cod\_turma' deveria ter ido para a tabela 'professores' como uma FK, mas no diagrama está ao contrário.
\choice Sim, pois as cardinalidades e os sentidos dos relacionamentos entre 'professores' e 'alunos' indica, no final, um relacionamento N:N entre essas duas tabelas.
\end{choices}

\question

\textit{Starvation} ocorre quando:
\begin{choices}
\choice Pelo menos um evento espera por um evento que não vai ocorrer.
\choice A prioridade de um processo é ajustada de acordo com o tempo total de execução do mesmo.
\choice O processo tenta mas não consegue acessar uma variável compartilhada.
\choice Pelo menos um processo é continuamente postergado e não executa.
\choice Dois ou mais processos são forçados a acessar dados críticos alternando estritamente entre eles.
\end{choices}

\question

Um sistema centralizado é um concentrador de recursos; um sistema distribuído apresenta seus recursos dispersos. Entretanto, nem todo conjunto de recursos computacionais dispersos pode ser considerado um sistema distribuído. Considerando um conjunto de computadores, assinale a alternativa que melhor corresponde às características necessárias para considerá-lo um sistema distribuído:
\begin{choices}
\choice Existência de sistema operacional idêntico e hardware padronizado em todos os computadores
\choice Suporte de rede e um relógio global
\choice Existência de memória compartilhada e relógios locais sincronizados 
\choice Existência de memória secundária compartilhada e protocolos de sincronização de estado
\choice Suporte de rede e funções primitivas de comunicação 
\end{choices}

\question

No programa abaixo, escrito em Pascal, os parâmetros do procedimento \texttt{vr} são passados por valor.



\begin{verbatim}

program teste;

var x, y: integer;



procedure vr(u, v: integer);

begin

    u := 2 * u;

    x := u + v;

    u := u - 1;

end;



begin

    x := 4;

    y := 2;

    vr(x, y);

    writeln(x);

end.

\end{verbatim}



O valor de \texttt{x} impresso na última linha do programa é:
\begin{choices}
\choice 8
\choice 4
\choice 5
\choice 10
\choice 7
\end{choices}

\question

Professor Mac Sperto propôs o seguinte algoritmo de ordenação, chamado de Super Merge, similar ao \textit{merge sort}: divida o vetor em 4 partes do mesmo tamanho (ao invés de 2, como é feito no \textit{merge sort}). Ordene recursivamente cada uma das partes e depois intercale-as por um procedimento semelhante ao procedimento de intercalação do \textit{merge sort}. Qual das alternativas abaixo é verdadeira?
\begin{choices}
\choice Super Merge está correto, mas consome tempo maior que \( O(\text{merge sort}) \)
\choice Super Merge não está correto. Não é possível ordenar quebrando o vetor em 4 partes
\choice Nenhuma das afirmações acima está correta
\choice Super Merge está correto, mas consome tempo \( O(\text{merge sort}) \)
\choice Super Merge está correto, mas consome tempo menor que \( O(\text{merge sort}) \)
\end{choices}

\question

Se \( \lim_{n \to \infty} \left( \frac{1}{n^2} + \frac{2}{n^2} + \cdots + \frac{n-1}{n^2} \right) = L \) então


\begin{choices}
\choice \( L = 1/2 \)
\choice \( L = 2 \)
\choice \( L = 1 \)
\choice \( L = 0 \)
\choice \( L = \infty \)
\end{choices}

\question

\begin{figure}[H]

    \centering

    \includegraphics[width=0.5\linewidth]{static/imgs/defaul.png}

    \caption{Enter Caption}

\end{figure}

Se \( O = (0, 0, 0) \); \( A = (2, 4, 1) \); \( B = (3, 1, 1) \) e \( C = (1, 3, 5) \) então o volume do sólido acima é


\begin{choices}
\choice 44
\choice 21
\choice 35/2
\choice 35
\choice 30
\end{choices}

\question

Qual o valor do atributo \( E.val \) após a análise da expressão "4 / 2 / 2" para o esquema de tradução a seguir?



\begin{itemize}

    \item \( E \rightarrow T / E1 \ \{ E.val = T.val / E1.val \} \)

    \item \( E \rightarrow T \ \{ E.val = T.val \} \)

    \item \( T \rightarrow \text{digito} \ \{ T.val = \text{val(digito)} \} \)

\end{itemize}


\begin{choices}
\choice 8
\choice 3
\choice 1
\choice 4
\choice 2
\end{choices}

\question

Sobre a técnica conhecida como Z-buffer é correto afirmar que:
\begin{choices}
\choice É possível realizar o cômputo das variáveis envolvidas de forma incremental.
\choice As primitivas geométricas precisam estar ordenadas de acordo com a distância em relação ao observador. 
\choice É uma técnica muito comum de detecção de colisão. 
\choice As dimensões do Z-buffer são independentes das dimensões do frame buffer. 
\choice Nenhuma das alternativas acima está correta. 
\end{choices}

\question

Suponha que \( T \) seja uma árvore AVL inicialmente vazia, e considere a inserção dos elementos \(10, 20, 30, 5, 15, 2\) em \( T \), nesta ordem. Qual das sequências abaixo corresponde a um percurso de \( T \) em pré-ordem:
\begin{choices}
\choice 15, 10, 5, 2, 20, 30
\choice 20, 10, 5, 2, 15, 30
\choice 10, 5, 2, 20, 15, 30
\choice 2, 5, 10, 15, 20, 30
\choice 30, 20, 15, 10, 5, 2
\end{choices}

\question

Uma integração de Sistemas Computacionais formando uma rede, tipicamente é implementada através da instalação de uma Arquitetura de Rede, que é composta de camadas e protocolos, em cada um dos elementos que compõem esta rede. Considere que estações “conversam” quando aplicações de usuários conseguem comunicar-se, sintática e semanticamente, através da Rede de Computadores. Baseados nesta premissa e em todos os conceitos associados à implementação e utilização das redes de computadores podemos afirmar como certo:


\begin{choices}
\choice Computadores com arquiteturas de redes diferentes conseguem “conversar”.
\choice Computadores com arquiteturas de redes diferentes podem “conversar” através de um \textit{gateway} ou conversor de protocolos.
\choice Computadores com arquiteturas diferentes podem “conversar” através de multiplexadores.
\choice Nenhuma delas é uma afirmação correta.
\choice Computadores com arquiteturas de rede parecidas conseguem “conversar”.
\end{choices}

\question

Considere o seguinte trecho de programa:\\

1. \( i := 1; \)\\

2. while \( i \leq n \) do\\

   begin\\

3. \( sum := sum + a[i]; \)\\

4. \( i := i + 1; \)\\

   end;\\



Considere que:\\

- \( I \) representa a inicialização da variável \( i := 1 \) na linha 1;\\

- \( T \) representa o teste da linha 2;\\

- \( A \) representa os comandos da linha 3;\\

- \( P \) representa o incremento na linha 4.\\



Qual é a expressão regular que representa todas as sequências de passos possíveis de serem executados por este trecho de programa?
\begin{choices}
\choice \( I(T(AP)^*)^+ \)
\choice \( I(TAP)^* \)
\choice \( IT^+A^*P^* \)
\choice \( I(TAP)^+ \)
\choice \( I(T(AP)^*)^* \)
\end{choices}

\question

Qual das seguintes opções \textbf{não é} uma vantagem do uso de um banco de

dados sobre o sistema de arquivos convencional?


\begin{choices}
\choice Eficiência
\choice Controle de acesso e visualização de dados
\choice Controle de redundância e inconsistência
\choice Independência de dados
\choice Estruturação de dados
\end{choices}

\question

Em um banco de dados relacional, o que é uma chave primária?


\begin{choices}
\choice É denominada por FK
\choice Nenhuma das respostas acima
\choice É um atributo (ou um conjunto de atributos) que permite identificar única e exclusivamente cada registro da tabela
\choice É uma validação que evita o cadastro de dados errados
\choice É qualquer coluna ou conjunto de colunas que podem servir como índice
\end{choices}

\question

Qual das informações a seguir \textbf{NÃO} é colocada no registro de ativação na chamada de funções?


\begin{choices}
\choice Variáveis locais estáticas
\choice Estado dos registradores
\choice Valor de retorno da função
\choice Endereço de retorno
\choice Link para a subrotina chamadora
\end{choices}

\question

Considerando a rede de Petri abaixo, quais das alternativas são verdadeiras?



\begin{itemize}

    \item[I)] O lugar \( A \) está habilitado a disparar.

    \item[II)] Apenas a transição \( T1 \) está habilitada a disparar.

    \item[III)] A sequência de transições \( (T1, T2, T3, T2) \) pode ser disparada, nessa ordem.

    \item[IV)] A transição \( T4 \) nunca poderá ser disparada.

\end{itemize}



\begin{figure}[H]

    \centering

    \includegraphics[width=0.5\linewidth]{static/imgs/defaul.png}

    \caption{Questão 49}

    \label{fig:enter-label}

\end{figure}
\begin{choices}
\choice Apenas as alternativas I e III.
\choice Apenas as alternativas III, IV.
\choice Todas as alternativas.
\choice Apenas as alternativas II, III e IV.
\choice Apenas as alternativas II e III.
\end{choices}

\question

Os protocolos de transporte atribuem a cada serviço um identificador único, o qual é empregado para encaminhar uma requisição de um aplicativo cliente ao processo servidor correto. Nos protocolos de transporte TCP e UDP, como esse identificador se denomina?
\begin{choices}
\choice Conexão
\choice Porta
\choice Protocolo de aplicação
\choice Identificador do processo (PID)
\choice Endereço IP
\end{choices}

\question

Uma árvore binária é declarada em C como



\begin{verbatim}

typedef struct no *apontador;

struct no {

    int valor;

    apontador esq, dir;

};

\end{verbatim}



onde \texttt{esq} e \texttt{dir} representam ligações para os filhos esquerdo e direito de um nó da árvore, respectivamente. Qual das seguintes alternativas é uma implementação correta da operação que inverte as posições dos filhos esquerdo e direito de um nó \texttt{p} da árvore, onde \texttt{t} é um apontador auxiliar.


\begin{choices}
\choice \texttt{t = p;} \\

          \texttt{p->esq = p->dir;} \\

          \texttt{p->dir = p->esq;}
\choice \texttt{p->esq = p->dir;} \\

          \texttt{t = p->esq;} \\

          \texttt{p->dir = t;}
\choice \texttt{t = p->dir;} \\

          \texttt{p->dir = p->esq;} \\

          \texttt{p->esq = t;}
\choice \texttt{t = p->dir;} \\

          \texttt{p->esq = p->dir;} \\

          \texttt{p->dir = t;}
\choice \texttt{p->dir = t;} \\

          \texttt{p->esq = p->dir;} \\

          \texttt{p->dir = t;}
\end{choices}

\question

Considere \( C(x) \) uma função que define a complexidade de um problema \( x \); \( E(x) \) uma função que define o esforço (em termos de tempo) exigido para se resolver o problema \( x \). Sejam dois problemas denominados \( p1 \) e \( p2 \). Assinale a alternativa correta.


\begin{choices}
\choice \( E(p1 + p2) < E(p1) + E(p2) \)
\choice \( C(p1 + p2) < C(p1) + C(p2) \)
\choice Se \( C(p1) < C(p2) \) então \( E(p1) < E(p2) \)
\choice Se \( C(p1) < C(p2) \) então \( E(p1) > E(p2) \)
\choice Nenhuma das alternativas anteriores
\end{choices}

\question

Você estava ouvindo a discussão de dois colegas, o João e a Maria. Enquanto o

João afirmava que o conceito de banco de dados (BD) é extremamente diferente do

conceito de sistema de gerenciamento de bancos de dados (SGBD), a Maria

discordava e afirmava que o banco de dados nada mais é do que um componente do

SGBD. Nesse momento o professor entrou na sala, percebeu a discussão e falou

que:
\begin{choices}
\choice Os dois estavam certos pois o conceito, na prática, depende da situação específica que se está considerando.
\choice Os dois estavam errados pois não são conceitos separados nem componentes, são conceitos sinônimos.
\choice O João estava certo já que o banco de dados é, realmente, um conceito diferente do conceito do sistema de gerenciamento de bancos de dados, e existe de forma independente.
\choice Não é possível saber quem estava certo ou errado pois não definiram qual o modelo de dados que estava sendo considerado na discusão.
\choice A Maria estava certa já que o banco de dados é, realmente, um componente interno do sistema de gerenciamento de bancos de dados e não existe de forma independente.
\end{choices}

\question

Assinale a alternativa INCORRETA:


\begin{choices}
\choice O controle de fluxo tem como objetivo garantir que nenhum dos parceiros de uma comunicação inunda o outro enviando pacotes mais rápido do que ele pode tratar.
\choice Os serviços orientados a conexão podem ajudar no controle de congestionamento através da diminuição da taxa de transmissão durante um congestionamento em andamento.
\choice Nos serviços orientados a conexões há a necessidade de estabelecimento de uma conexão antes da transferência dos dados.
\choice Serviços orientados a conexão podem ser implementados em subredes que funcionam no modo datagrama.
\choice Os serviços orientados a conexões são sempre confiáveis garantindo a entrega ordenada e completa dos dados transmitidos.
\end{choices}

\question

Analise o diagrama relacional do banco de dados exibido na figura abaixo:

\begin{figure}[H]

  \centering

  \caption{Banco de dados de peças e fornecedores}

  \label{fig:pecas}

  %\fbox{

    \includegraphics[width=0.5\linewidth]{static/imgs/defaul.png}

  %}\\

  %\footnotesize{Fonte: xxxx.}

\end{figure}

Podemos dizer que esse banco de dados ilustra:
\begin{choices}
\choice Um relacionamento com cardinalidade N:N entre as tabelas ``pecas' e ``fornecimentos', realizado através de relacionamentos não identificados com a tabela ``fornecedores'. 
\choice Um relacionamento com cardinalidade N:N entre as tabelas ``pecas' e ``fornecedores', realizado através de relacionamentos identificados através da tabela ``fornecimentos'. 
\choice Um relacionamento com cardinalidade N:N entre as tabelas ``pecas' e ``fornecedores', realizado através de relacionamentos não identificados através da tabela ``fornecimentos'. 
\choice Um relacionamento com identificação N:N entre as tabelas ``pecas' e ``fornecedores', realizado através do relacionamento com cardinalidade identificada através da tabela ``fornecimentos'. 
\choice Um relacionamento com cardinalidade 1:N entre as tabelas ``pecas' e ``fornecedores', indicando que uma peça pode ter sido fornecida por diversos fornecedores. 
\end{choices}

\question

Considere as seguintes afirmações sobre resolução de problemas em IA.

\begin{enumerate}[label=\Roman*.]

    \item A* é um conhecido algoritmo de busca heurística.

    \item O Minimax é um dos principais algoritmos para jogos de dois jogadores, como o xadrez.

    \item Busca em espaço de estados é uma das formas de resolução de problemas em IA.

\end{enumerate}

São corretas:
\begin{choices}
\choice Apenas III
\choice Apenas I e II
\choice Apenas I e III
\choice Apenas II e III
\choice I, II e III
\end{choices}

\question

Dado um vetor \( u \in \mathbb{R}^2 \), \( u = (-3, 4) \), vamos denotar por \( v \) o vetor de \( \mathbb{R}^2 \) que tem tamanho 1 e é ortogonal a \( u \). Então \( v \) pode ser dado por


\begin{choices}
\choice \( \left( -\frac{4}{5}, \frac{3}{5} \right) \)
\choice \( \left( -\frac{4}{5}, -\frac{3}{5} \right) \)
\choice \( \left( -\frac{4}{5}, \frac{2}{5} \right) \)
\choice \( \left( -\frac{4}{5}, \frac{1}{5} \right) \)
\choice \( \left( \frac{3}{5}, \frac{4}{5} \right) \)
\end{choices}

\question

Assinale a proposição verdadeira


\begin{choices}
\choice Se \( x \) e \( y \) são números reais tais que \( x < y \) então \( x^2 < y^2 \)
\choice Se \( x \) é um número real tal que \( x^2 \leq 4 \) então \( x \leq 2 \) e \( x \leq -2 \)
\choice nenhuma das anteriores
\choice Se \( x < -4 \) ou \( x > 1 \) então \( \frac{2x + 3}{x - 1} > 1 \)
\choice Se \( x + y \) é um número racional então \( x \) e \( y \) são números racionais
\end{choices}

\question

(ENADE 2017) No modelo relacional de dados, uma junção entre tabelas cria uma

outra tabela derivada de duas ou mais tabelas de acordo com os critérios

especificados para a junção, de modo similar às regras da teoria dos

conjuntos. Nos conjuntos a seguir, as áreas hachuradas representam os resultados

das consultas SQL em que ``id' é uma chave primária e ``fk' é uma chave

estrangeira:

\begin{figure}[H]

  \centering

  \includegraphics[width=0.5\linewidth]{static/imgs/defaul.png}

\end{figure}

Assinale a opção em que a declaração SQL é equivalente ao conjunto apresentado a

seguir:

\begin{figure}[H]

  \centering

  \includegraphics[width=0.5\linewidth]{static/imgs/defaul.png}

\end{figure}


\begin{choices}
\choice \begin{verbatim}SELECT *

FROM tabela1 t1

WHERE NOT EXISTS (SELECT 1

                  FROM tabela1 t2

                  WHERE t2.id = tk.fk);

\end{verbatim}


\choice \begin{verbatim}SELECT *

FROM tabela1 t1 WHERE EXISTS (SELECT 1

                              FROM tabela2 t2

                              WHERE t2.id = t1.fk);

\end{verbatim}


\choice \begin{verbatim}SELECT *

FROM tabela1 t1

RIGHT JOIN tabela2 t2 ON (t1.id = t2.fk)

WHERE t1.id IS NULL;

\end{verbatim}


\choice \begin{verbatim}SELECT *

FROM tabela1 t1

RIGHT JOIN tabela2 t2 ON (t1.id = t2.fk)

WHERE t2.fk <> NULL;

\end{verbatim}


\choice \begin{verbatim}SELECT *

FROM tabela1 t1

INNER JOIN tabela2 t2 ON (t1.id = t2.fk);

\end{verbatim}


\end{choices}

\question

O processo de visualização de objetos 3D envolve uma série de passos desde a representação vetorial de um objeto até a exibição da imagem correspondente na tela do computador (pipeline 3D). Selecione a alternativa abaixo que reflete a ordem correta em que esses passos devem ocorrer. 
\begin{choices}
\choice Recorte 3D, transformação de câmera, rasterização, projeção, mapeamento para coordenadas de tela. 
\choice Nenhuma das respostas acima está correta.
\choice Transformação de câmera, mapeamento para coordenadas de tela, recorte 3D, rasterização, projeção. 
\choice Transformação de câmera, recorte 3D, projeção, mapeamento para coordenadas de tela, rasterização. 
\choice Projeção, transformação de câmera, recorte 3D, mapeamento para coordenadas de tela, rasterização. 
\end{choices}

\question

Considere as seguintes afirmações sobre mecanismos de inferência em sistemas baseados em regras.

\begin{enumerate}[label=\Roman*.]

    \item O encadeamento regressivo tem pouca utilidade prática, pois deve partir do possível resultado.

    \item O encadeamento progressivo tanto pode ser em amplitude quanto em profundidade.

    \item Podem trabalhar com informações incertas ou incompletas. 

\end{enumerate}

São corretas:
\begin{choices}
\choice Apenas II e III
\choice Apenas I e II
\choice Apenas I e III
\choice Apenas III
\choice I, II e III
\end{choices}

\question

Todos os convidados presentes num jantar tomam chá ou café. Treze convidados bebem café, dez bebem chá e 4 bebem chá e café. Quantas pessoas tem nesse jantar.


\begin{choices}
\choice 10
\choice 19
\choice 27
\choice 23
\choice 15
\end{choices}

\question

Quais dos algoritmos de ordenação abaixo possuem tempo no pior caso e tempo médio de execução proporcional a \( O(n \log n) \).
\begin{choices}
\choice Heap sort e selection sort
\choice Quicksort e merge sort
\choice Merge sort e bubble sort
\choice Bubble sort e Quick sort
\choice Merge sort e heap sort
\end{choices}

\question

Supondo a Relação $\texttt{PROJ (PNO, Orçam)}$, com chave primária $\texttt{PNO}$, a Relação $\texttt{EMP (ENO, ENome, Cargo)}$ com chave primária $\texttt{ENO}$, e a Relação $\texttt{DSG (ENO, PNO, Dur, Resp)}$, com chave primária \{\texttt{ENO, PNO}\}, chave estrangeira $\texttt{PNO}$ em relação a $\texttt{PROJ}$ e chave estrangeira $\texttt{ENO}$ em relação a $\texttt{EMP}$. Qual das expressões da álgebra relacional abaixo $\textbf{NÃO}$ corresponde à seguinte consulta SQL:



\begin{verbatim}

SELECT ENome

FROM EMP, PROJ, DSG

WHERE EMP.ENO = DSG.ENO

      AND PROJ.PNO = DSG.PNO

      AND Dur > 36

\end{verbatim}
\begin{choices}
\choice \( \pi_{ENome} (\text{PROJ} \Join_{PNO} (\text{EMP} \Join_{ENO} \sigma_{Dur > 36} (\pi_{Dur} (\text{DSG})))) \)
\choice \( \pi_{ENome} (\text{PROJ} \Join_{PNO} (\text{EMP} \Join_{ENO} \sigma_{Dur > 36} (\text{DSG}))) \)
\choice \( \pi_{ENome} (\sigma_{Dur > 36} ((\pi_{PNO} (\text{PROJ})) \Join_{PNO} (\text{EMP} \Join_{ENO} (\pi_{Dur} (\text{DSG})))) \)
\choice \( \pi_{ENome} (\text{PROJ} \Join_{PNO} (\sigma_{Dur > 36} (\text{EMP} \Join_{ENO} (\text{DSG})))) \)
\choice \( \pi_{ENome} (\sigma_{Dur > 36} ((\pi_{ENome, ENO} (\text{EMP})) \Join_{ENO} (\text{PROJ} \Join_{PNO} (\text{DSG})))) \)
\end{choices}

\question

A função abaixo computa a soma dos \( n \) primeiros números inteiros não negativos:



\begin{verbatim}

function sum(n: integer): integer;

begin

    if n = 0 then sum := 0

    else -------------

end;

\end{verbatim}



A parte que falta para completar a condição \texttt{else} é:
\begin{choices}
\choice \texttt{while n<>0 sum := sum + sum(n+1)}
\choice \texttt{sum := (n-1) + sum(n)}
\choice \texttt{sum := (n-1) + sum(n-1)}
\choice \texttt{sum := n + sum(n)}
\choice \texttt{sum := n + sum(n-1)}
\end{choices}

\question

Quando trabalhando com sistemas baseados em trocas de mensagens, temporizações (\textit{time-outs}) são utilizadas para:


\begin{choices}
\choice Temporariamente suspender a transmissão de mensagens.
\choice Arbitrar que uma mensagem transmitida foi perdida.
\choice Limitar o número de retransmissões de uma mensagem.
\choice Limitar o tempo para obter um recurso.
\choice Limitar o tamanho de uma mensagem transmitida.
\end{choices}

\question

Ao executar o seguinte comando SQL, nenhum dado foi retornado, ou seja, não

ocorreu nenhum erro de sintaxe mas nenhum dado foi encontrado:

\begin{tcolorbox}

 \begin{verbatim}1   SELECT c.nome AS cliente

 2   FROM clientes c

 3   LEFT JOIN pedidos p ON (p.cliente_id = c.cliente_id)

 4   WHERE p.cliente_id IS NULL

 5   ORDER BY cliente;

 \end{verbatim}

\end{tcolorbox}

Por que esse comando SQL não retornou nenhum dado?


\begin{choices}
\choice Há um erro semântico na linha 3
\choice Há um erro lógico na linha 4.
\choice A junção que deveria ter sido realizada, na linha 3, era uma FULL JOIN ao invés de uma RIGHT JOIN.
\choice Todos os clientes cadastrados fizeram, pelo menos, um pedido de compra.
\choice A junção que deveria ter sido realizada, na linha 3, era uma RIGHT JOIN ao invés de uma LEFT JOIN.
\end{choices}

\question

O Modelo Espiral de desenvolvimento de software, proposto por Barry Boehm na década de 1970, apresenta as seguintes afirmações:

\begin{enumerate}[label=\Roman*.]

    \item Suas atividades de desenvolvimento são conduzidas por riscos;

    \item Cada ciclo da espiral inclui 4 passos: passo 1 - identificação dos objetivos; passo 2 - avaliação das alternativas tendo em vista os objetivos e os riscos (incertezas, restrições) do desenvolvimento; passo 3 - desenvolvimento de estratégias (simulação, prototipagem) para resolver riscos; e passo 4 - planejamento do próximo passo e continuidade do processo determinada pelos riscos restantes;

    \item É um modelo evolutivo em que cada passo pode ser representado por um quadrante num diagrama cartesiano: assim, na dimensão radial da espiral tem-se o custo acumulado dos vários passos do desenvolvimento, enquanto na dimensão angular tem-se o progresso do projeto.

\end{enumerate}

Levando-se em conta as três afirmações I, II e III acima, identifique a única alternativa válida:


\begin{choices}
\choice Apenas a I e a III estão corretas;
\choice Apenas a I e a II estão corretas;
\choice Apenas a III está correta.
\choice Apenas a II e a III estão corretas;
\choice As afirmações I, II e III estão corretas;
\end{choices}

\question

Qual das seguintes afirmações sobre crescimento assintótico de funções não é verdadeira:
\begin{choices}
\choice Se \( f(n) = O(g(n)) \) então \( g(n) = O(f(n)) \)
\choice \( 2^{n+1} = O(2^n) \)
\choice \( 2n^2 + 3n + 1 = O(n^2) \)
\choice Se \( f(n) = O(g(n)) \) e \( g(n) = O(h(n)) \) então \( f(n) = O(h(n)) \)
\choice \( \log n^2 = O(\log n) \)
\end{choices}

\question

Sobre a UML, quais das seguintes afirmações são verdadeiras?



\begin{itemize}

    \item[I)] A UML é o método de desenvolvimento de software mais utilizado na atualidade.

    \item[II)] A UML é uma evolução das linguagens para especificação dos conceitos dos métodos de Booch, OMT e OOSE e também de outros métodos de especificação de requisitos de software orientados a objetos ou não.

    \item[III)] A UML é composta dos seguintes diagramas: Diagrama de Caso de Uso, Diagrama de Classes, Diagrama de Colaboração, Diagrama de Estados, entre outros.

    \item[IV)] Em UML pode-se representar tão somente relacionamentos de Agregação, Associação e Composição.

\end{itemize}
\begin{choices}
\choice Apenas as alternativas III e IV.
\choice Apenas as alternativas II e III.
\choice Apenas as alternativas I, II e III.
\choice Todas as alternativas.
\choice Nenhuma delas.
\end{choices}

\question

Na Álgebra Booleana:
\begin{choices}
\choice Os dígitos são hexadecimais de 0 a 15
\choice Os dígitos são octais, de 0 a 7
\choice Há dez valores motivados pelos dez dedos do ser humano
\choice Os dígitos são binários 0 e 1
\choice Os dígitos são alfanuméricos que podem ser representados por um byte
\end{choices}

\question

Considere as seguintes afirmativas sobre as linguagens usadas para análise sintática:

\begin{enumerate}[label=\Roman*.]

    \item A classe LL(1) não aceita linguagens com produções que apresentem recursões diretas à esquerda (ex. \(L \rightarrow La\)) mas aceita linguagens com recursões indiretas (ex. \(L \rightarrow Ra, R \rightarrow Lb\)).

    \item A linguagem LR(1) reconhece a mesma classe de linguagens que LALR(1).

    \item A linguagem SLR(1) reconhece uma classe de linguagens maior que LR(0).

\end{enumerate}

Selecione a afirmativa correta:
\begin{choices}
\choice As afirmativas II e III são verdadeiras
\choice Apenas a afirmativa III é verdadeira
\choice As afirmativas I e II são verdadeiras
\choice As afirmativas I e III são falsas
\choice As afirmativas I e III são verdadeiras
\end{choices}

\question

Quando Edgar Codd cricou o modelo relacional de dados,

também criou algumas regras para que a estrutura das tabelas fosse adequada às

operações de inserção, atualização, remoção e busca de dados. Ao avaliar um

banco de dados que armazena dados dos eleitores, dos partidos e dos candidatos

de uma eleição, você encontrou a situação ilustrada na tabela abaixo:

\begin{figure}[H]

  \centering

  \includegraphics[width=0.5\linewidth]{static/imgs/defaul.png}

\end{figure}

Em relação às boas normas de projeto do modelo relacional, assinale a

alternativa que corretamente identifica um problema no projeto desta tabela:


\begin{choices}
\choice Apesar do telefone ser um atributo multivalorado, do jeito que a tabela está criada só é possível armazenar dois números de telefones, criando uma restrição ao armazenamento de telefones de contado de eleitores que tenham mais de dois telefones.
\choice A tabela contém um atributo multivalorado (o correto seria que cada tabela não tivesse atributos multivalorados).
\choice A chave primária para a tabela não está definida e, somente olhando as colunas, não é possível dizer se é a coluna 'numero' ou a coluna 'titulo'.
\choice Todas as afirmativas anteriores.
\choice A tabela está misturando dados de três entidades diferentes, eleitores, partidos e candidatos (o correto seria que cada tabela armazenasse dados de um único assunto ou entidade).
\end{choices}

\question

Considere um grafo \( G \) satisfazendo as seguintes propriedades:



\begin{itemize}

    \item \( G \) é conexo

    \item Se removermos qualquer aresta de \( G \), o grafo obtido é desconexo.

\end{itemize}



Então é correto afirmar que o grafo \( G \) é:
\begin{choices}
\choice Uma árvore
\choice Hamiltoniano
\choice Um circuito
\choice Não bipartido
\choice Euleriano
\end{choices}

\question

Em relação ao paradigma de programação cliente-servidor. Qual das afirmativas abaixo é FALSA?
\begin{choices}
\choice Um aplicativo cliente é um programa arbitrário que se torna temporariamente um cliente quando for necessário o acesso remoto a um serviço, mas também executa processamento local.
\choice Um aplicativo servidor é um programa de propósito especial dedicado a fornecer um serviço, mas pode tratar de múltiplos clientes remotos ao mesmo tempo. 
\choice Um aplicativo cliente pode acessar múltiplos serviços quando necessário.
\choice Um aplicativo servidor inicia ativamente o contato com clientes 
\choice Um aplicativo servidor aceita contato de clientes arbitrários, mas oferece um único serviço. 
\end{choices}

\question

Histograma de uma imagem com \( K \) tons de cinza é:
\begin{choices}
\choice Contagem do número de vezes que cada um dos \( K \) tons de cinza ocorreu na imagem.
\choice Contagem do número de tons de cinza que ocorreram na imagem.
\choice Contagem do número de objetos encontrados na imagem.
\choice Contagem dos pixels da imagem.
\choice Nenhuma alternativa acima.
\end{choices}



    %%%%%%%%%%%%%%%%%%%%%%%%%%%%%%%%%%%%%%%%%%%%%%%%%%%%%%%%%%%%%%%%%%%%%%%%%%%%%%%%
    %%% Finaliza o bloco de questões:
    \end{questions}
    \end{document}
    